\documentclass[12pt]{report}

%% ─── USER-CONFIGURABLE VARIABLES ────────────────────────────────────────────
\newcommand{\ProjectTitle}{Flood Sentinel: Self-Supervised Temporal Representation Learning for Early Flood Warning from USGS Streamgage Networks}
\newcommand{\CourseCode}{COMP 493}
\newcommand{\YearSemester}{Fourth Year / Second Semester}
\newcommand{\StudentOne}{Aditya Shrestha}
\newcommand{\RollNo}{Roll No. 57}
\newcommand{\SupervisorName}{Pankaj Raj Dawadi}
\newcommand{\SupervisorDesignation}{Associate Professor}
\newcommand{\SupervisorDept}{Department of Computer Science and Engineering, School of Engineering}
\newcommand{\SubmissionDate}{21 September 2026}
%% ─────────────────────────────────────────────────────────────────────────────

%% Fonts - Times New Roman equivalent
\usepackage{newtxtext, newtxmath}

%% Page geometry - exact KU DoCSE margins
\usepackage[
  letterpaper,
  left=1.5in,
  right=1.25in,
  top=1.25in,
  bottom=1.25in,
  headheight=14pt,
  headsep=0pt,
  footskip=0.5in
]{geometry}

%% Line and paragraph spacing
\usepackage{setspace}
\setstretch{1.5}
\setlength{\parindent}{0pt}
\setlength{\parskip}{18pt}

%% Headers and footers
\usepackage{fancyhdr}
\pagestyle{fancy}
\fancyhf{}
\fancyfoot[C]{\thepage}
\renewcommand{\headrulewidth}{0pt}
\renewcommand{\footrulewidth}{0pt}

\fancypagestyle{plain}{%
  \fancyhf{}
  \fancyfoot[C]{\thepage}
  \renewcommand{\headrulewidth}{0pt}
}

%% Heading formatting
\usepackage{titlesec}

\titleformat{\chapter}[hang]
  {\fontsize{16pt}{19.2pt}\selectfont\bfseries\setstretch{1.0}}
  {Chapter~\thechapter}
  {1.5em}
  {}
\titlespacing*{\chapter}{0pt}{0pt}{18pt}

\titleformat{\section}[hang]
  {\fontsize{14pt}{16.8pt}\selectfont\bfseries\setstretch{1.0}}
  {\thesection}
  {1em}
  {}
\titlespacing*{\section}{0pt}{18pt}{0pt}

\titleformat{\subsection}[hang]
  {\fontsize{13pt}{15.6pt}\selectfont\bfseries\setstretch{1.0}}
  {\thesubsection}
  {1em}
  {}
\titlespacing*{\subsection}{1.5em}{10pt}{0pt}

\titleformat{\subsubsection}[hang]
  {\fontsize{12pt}{14.4pt}\selectfont\bfseries\setstretch{1.0}}
  {\thesubsubsection}
  {1em}
  {}
\titlespacing*{\subsubsection}{3em}{10pt}{0pt}

%% Table of Contents formatting
\usepackage{titletoc}
\usepackage{tocloft}

\renewcommand{\cfttoctitlefont}{\hfill\fontsize{16pt}{19.2pt}\selectfont\bfseries}
\renewcommand{\cftaftertoctitle}{\hfill}
\setlength{\cftbeforetoctitleskip}{0pt}
\setlength{\cftaftertoctitleskip}{18pt}

\renewcommand{\cftchappresnum}{Chapter~}
\renewcommand{\cftchapaftersnum}{}
\setlength{\cftchapnumwidth}{6em}
\renewcommand{\cftchapleader}{\cftdotfill{\cftdotsep}}
\renewcommand{\cftchapfont}{}
\renewcommand{\cftchappagefont}{}

\setlength{\cftsecindent}{6em}
\setlength{\cftsecnumwidth}{2.5em}
\setlength{\cftsubsecindent}{9em}
\setlength{\cftsubsecnumwidth}{3em}

\setcounter{tocdepth}{3}
\setcounter{secnumdepth}{4}

\renewcommand{\cftloftitlefont}{\hfill\fontsize{16pt}{19.2pt}\selectfont\bfseries}
\renewcommand{\cftafterloftitle}{\hfill}
\renewcommand{\cftlottitlefont}{\hfill\fontsize{16pt}{19.2pt}\selectfont\bfseries}
\renewcommand{\cftafterlottitle}{\hfill}

%% Caption formatting
\usepackage{caption}
\captionsetup{
  font={bf, footnotesize},
  labelsep=space,
  justification=centering,
  singlelinecheck=false
}
\captionsetup[table]{position=above}
\captionsetup[figure]{position=below}

\usepackage{chngcntr}
\counterwithin{figure}{chapter}
\counterwithin{table}{chapter}

%% Graphics and utilities
\usepackage{graphicx}
\usepackage[style=apa, backend=biber]{biblatex}
\addbibresource{references.bib}
\usepackage{microtype}
\usepackage{enumitem}
\usepackage{booktabs}
\usepackage{array}
\usepackage{tikz}
\usetikzlibrary{arrows.meta, positioning, shapes.geometric}
\usepackage{hyperref}
\hypersetup{colorlinks=false, hidelinks}


\begin{document}

%% ─── TITLE PAGE ─────────────────────────────────────────────────────────────
\begin{titlepage}
  \centering
  \setstretch{1.0}
  \setlength{\parskip}{0pt}

  {\fontsize{14pt}{16.8pt}\selectfont\bfseries Kathmandu University\par}
  {\fontsize{14pt}{16.8pt}\selectfont\bfseries Department of Computer Science and Engineering\par}
  {\fontsize{14pt}{16.8pt}\selectfont\bfseries Dhulikhel, Kavre\par}

  \vspace{1cm}

  \includegraphics[width=1.2in, height=1.2in, keepaspectratio]{ku_logo.png}

  \vspace{0.8cm}

  {\fontsize{12pt}{14.4pt}\selectfont\bfseries A Project Proposal\par}
  {\fontsize{12pt}{14.4pt}\selectfont\bfseries on\par}
  {\fontsize{12pt}{14.4pt}\selectfont\bfseries ``\ProjectTitle''\par}

  \vspace{0.5cm}

  {\fontsize{12pt}{14.4pt}\selectfont\bfseries [Code No: \CourseCode]\par}
  {\fontsize{12pt}{14.4pt}\selectfont\bfseries
    (For partial fulfillment of \YearSemester{} in Computer Science and Engineering)\par}

  \vspace{1.2cm}

  {\fontsize{12pt}{14.4pt}\selectfont\bfseries Submitted by\par}

  \vspace{0.3cm}

  {\fontsize{12pt}{14.4pt}\selectfont\bfseries \StudentOne\par}
  {\fontsize{12pt}{14.4pt}\selectfont\bfseries \RollNo\par}
  \vspace{1.2cm}

  {\fontsize{12pt}{14.4pt}\selectfont\bfseries Submitted to\par}

  \vspace{0.3cm}

  {\fontsize{12pt}{14.4pt}\selectfont\bfseries \SupervisorName\par}
  {\fontsize{12pt}{14.4pt}\selectfont\bfseries Department of Computer Science and Engineering\par}

  \vspace{1cm}

  {\fontsize{12pt}{14.4pt}\selectfont\bfseries Submission Date: \SubmissionDate\par}

\end{titlepage}
%% ─────────────────────────────────────────────────────────────────────────────

%% ─── BONA FIDE CERTIFICATE ──────────────────────────────────────────────────
\thispagestyle{empty}
\setstretch{1.0}
\setlength{\parskip}{0pt}

\begin{center}
  {\fontsize{16pt}{19.2pt}\selectfont\bfseries Bona fide Certificate\par}
\end{center}

\vspace{3cm}

\begin{center}
  {\bfseries This project proposal on\par}
  \vspace{0.3cm}
  {\bfseries ``\ProjectTitle''\par}
  \vspace{0.3cm}
  {\bfseries is the bona fide work of\par}
  \vspace{0.3cm}
  {\bfseries ``\StudentOne'' (\RollNo)\par}
  \vspace{0.3cm}
  {\bfseries who prepared the project proposal under my supervision.\par}
\end{center}

\vspace{4cm}

{\bfseries Project Supervisor\par}
\vspace{1.5cm}
{\bfseries \rule{6cm}{0.4pt}\par}
{\bfseries \SupervisorName\par}
{\bfseries \SupervisorDesignation\par}
{\bfseries \SupervisorDept\par}
\vspace{0.8cm}
{\bfseries Date:\par}

\newpage
\setstretch{1.5}
\setlength{\parskip}{18pt}
%% ─────────────────────────────────────────────────────────────────────────────

%% ─── FRONT MATTER ───────────────────────────────────────────────────────────
\pagenumbering{roman}

\chapter*{Abstract}
\addcontentsline{toc}{chapter}{Abstract}

Public streamgage networks hold decades of daily discharge and stage records, but flood-event labels are scarce. This project tests whether a causal encoder pretrained on those unlabeled records by masked reconstruction produces anomaly scores that rise before a gauge crosses its official NWS flood stage. The encoder combines dilated causal convolutions with a two-layer causal transformer. It is trained on daily discharge, gage height, gridMET meteorology, and SNODAS snow water equivalent for 54 streamgages across 14 HUC-2 regions in the contiguous United States. Scores are compared with four baseline families: the National Water Model v3.0 retrospective simulation; persistence, climatology, and changepoint detectors; Isolation Forest and an LSTM autoencoder; and an Entity-Aware LSTM trained with 1\% to 100\% of the event labels. Five hypotheses and their stopping rules are pre-specified in this document before any test-gauge result is computed. The study is retrospective, uses daily data, and does not address flash floods.

\bigskip
\noindent\textbf{Keywords:} \textit{self-supervised learning, flood early warning, anomaly detection, USGS streamgage, temporal convolutional network, transformer, hydrometric time series}
\newpage
\tableofcontents


\chapter*{List of Figures}
\addcontentsline{toc}{chapter}{List of Figures}
\listoffigures
\newpage

\chapter*{List of Tables}
\addcontentsline{toc}{chapter}{List of Tables}
\listoftables
\newpage

\chapter*{Acronyms/Abbreviations}
\addcontentsline{toc}{chapter}{Acronyms/Abbreviations}
\begin{tabbing}
  \hspace{3cm} \= \kill
  \textbf{AUC-PR} \> Area Under the Precision-Recall Curve \\
  \textbf{AUC-ROC}\> Area Under the Receiver Operating Characteristic Curve \\
  \textbf{CAMELS} \> Catchment Attributes and Meteorology for Large-sample Studies \\
  \textbf{CONUS}  \> Contiguous United States \\
  \textbf{CUSUM}  \> Cumulative Sum control chart \\
  \textbf{EA-LSTM}\> Entity-Aware Long Short-Term Memory \\
  \textbf{EWMA}   \> Exponentially Weighted Moving Average \\
  \textbf{GAGES}  \> Geospatial Attributes of Gages for Evaluating Streamflow \\
  \textbf{HAC}    \> Heteroskedasticity and Autocorrelation Consistent \\
  \textbf{HUC}    \> Hydrologic Unit Code \\
  \textbf{NOAA}   \> National Oceanic and Atmospheric Administration \\
  \textbf{NWM}    \> National Water Model \\
  \textbf{NWPS}   \> National Water Prediction Service \\
  \textbf{SNODAS} \> Snow Data Assimilation System \\
  \textbf{SSL}    \> Self-Supervised Learning \\
  \textbf{TCN}    \> Temporal Convolutional Network \\
  \textbf{USGS}   \> United States Geological Survey \\
\end{tabbing}
\newpage
%% ─────────────────────────────────────────────────────────────────────────────

%% ─── MAIN MATTER ─────────────────────────────────────────────────────────────
\pagenumbering{arabic}
\setcounter{page}{1}

\chapter{Introduction}

Streamgages record discharge and river stage daily for decades, yet almost none of that record carries flood event labels. In contrast, flood catalogs are small, assembled by hand, and irregular across monitoring jurisdictions. Between 2014 and 2023, freshwater flooding caused an average of roughly 100 fatalities per year in the United States (National Weather Service, 2024). Operational hydrological forecasting relies on numerical simulations such as the National Water Model, but issuing timely alerts on medium-to-large catchments remains difficult due to regional calibration errors and high computational demands.

Data-driven hydrological models demonstrate that neural networks can match or outperform physical simulations for streamflow forecasting. However, existing machine learning workflows almost universally require supervised training labels. This reliance restricts model development to gauges with curated flood records, which are scarce and time-consuming to compile.

Self-supervised learning offers an alternative: temporal encoders can train directly on continuous, unlabeled streamflow observations without requiring event annotations. This proposal presents \textit{Flood Sentinel}. I investigate whether self-supervised pretraining on multi-channel streamgage time series produces anomaly scores that rise before a gauge exceeds its official flood stage, and whether those scores provide predictive utility beyond empirical baselines and operational simulations.

\section{Background}

The United States Geological Survey (USGS) operates over 10,000 streamgages across the contiguous United States, many providing continuous records of discharge and stage spanning 20 to 50 years. Combining these daily hydrometric observations with gridded surface meteorology from gridMET (Abatzoglou, 2013), snow water equivalent estimates from SNODAS, and physical catchment attributes from GAGES-II (Falcone, 2011) and CAMELS (Addor et al., 2017) produces an observational record across diverse hydro-climatic regimes.

Recent studies explore self-supervised learning for hydrological time series, but their objectives differ from flood early warning. HydroGEM (Haq et al., 2025) applied masked pretraining to streamflow records specifically for sensor quality control and data cleaning, not event detection. Global foundation models such as AIFL (Taccari et al., 2026) train multi-million parameter networks across thousands of catchments, requiring high-performance computing clusters that lie outside standard academic workstation setups. Other frameworks, such as Errorcastnet (Tran et al., 2025), focus on learning residual error corrections for physical models rather than detecting precursor anomaly signals directly from sensor observations.

Whether a causal temporal encoder trained exclusively on unlabeled hydrometric series provides actionable warning lead times before river stage crosses flood thresholds remains untested. Furthermore, prior machine learning applications in hydrology have frequently suffered from subtle temporal leakage and improper spatial validation splits (Kapoor \& Narayanan, 2023). This project addresses both issues by establishing publication-time filtering rules and frozen spatial-temporal partitions.

\section{Objectives}

\begin{itemize}
  \item To develop a causal temporal encoder combining dilated convolutions and transformer self-attention layers that learns representations from continuous streamgage observations without event annotations.
  \item To test whether anomaly scores from these representations detect impending flood events across holdout gauges earlier than empirical, physical, and supervised baselines.
  \item To implement a leakage-free data processing pipeline that applies publication-time filtering and strict spatial-temporal holdout splits.
\end{itemize}

\section{Motivation and Significance}

Flood early warning provides tangible value when emergency managers receive reliable advance notice. Even an additional 24 to 48 hours of warning allows communities to position temporary barriers, clear drainage channels, and coordinate evacuations. While self-supervised representation learning is standard practice in computer vision and natural language processing, its utility for continuous hydrometric series remains under-evaluated. The multi-decade records maintained by national hydrometric networks provide an empirical testbed for these methods.

Prior data-driven studies in hydrology have faced three specific difficulties that I address in this design. First, instead of relying on manually curated flood catalogs, I pretrain the encoder on raw time series without event annotations. Second, I guard against train-test contamination by enforcing publication-time filters and spatial-temporal holdout partitions. Third, rather than comparing against a single baseline, I evaluate against four distinct baseline families: physical simulation (NWM v3.0 Retrospective), empirical statistical filters, unsupervised anomaly detectors, and a supervised EA-LSTM ablated across six label budgets.

This project does not compete with the NWM or aim at a continental foundation model. It asks one question: given only unlabeled gauge records, does a causal encoder produce a signal that rises before a gauge crosses flood stage, and does that beat persistence, climatology, and a supervised EA-LSTM at small label budgets? Chapter~4 lists the outcomes I have committed to reporting, whatever they are.


\chapter{Related Works}

\section{Self-Supervised Learning for Time Series}

Self-supervised learning enables representation learning from unlabeled sequential observations. TS2Vec (Yue et al., 2022) applied hierarchical contrastive learning across temporal slices and instance sub-series, while TNC (Tonekaboni et al., 2021) used stationary temporal neighborhoods to build invariant representations. PatchTST (Nie et al., 2023) demonstrated that masked autoencoding across sub-series patches yields competitive forecasting representations for multivariate time series.

However, standard time series masked autoencoders often employ bidirectional attention, where masked values attend to both past and future timesteps. In an operational hydrological setting, attending to future observations introduces lookahead leakage. With causal convolutions and causal attention masks, a masked timestep can only be reconstructed from earlier context. The pretext task in this work is therefore framed as causal autoregressive masked reconstruction, not bidirectional masked autoencoding. I considered contrastive pretraining (such as TS2Vec and TNC) but rejected it: contrastive objectives require designing domain-specific temporal augmentations that can alter physical mass-balance dynamics, whereas masked reconstruction preserves the physical relationships across input channels.

\section{Data-Driven Hydrological Modeling}

Machine learning for rainfall-runoff modeling has expanded rapidly over the past decade. Kratzert et al. (2019) introduced the Entity-Aware LSTM (EA-LSTM) for regional streamflow modeling using CAMELS catchment attributes, showing that LSTMs generalize across diverse catchments when conditioned on static soil and topographic features. I adopt their open-source EA-LSTM architecture as the supervised reference baseline. In transformer modeling, Liu et al. (2022) demonstrated the utility of self-attention encoders for monthly streamflow predictions on the Yangtze River. Evaluating deep learning models across the Mekong River Basin, Nguyen et al. (2023) found that recurrent architectures frequently outperformed self-attention networks at extended multi-day forecast horizons. Because my encoder incorporates self-attention layers, Nguyen et al.'s findings directly motivate an ablation study: I include a TCN-only variant alongside the full TCN-Transformer to test whether self-attention adds measurable benefit over dilated convolutions alone.

In self-supervised hydrological modeling, HydroGEM (Haq et al., 2025) pretrained a 14.2M-parameter encoder on 3,724 USGS streamgages for sensor quality control and anomaly imputation. However, its model weights and pretraining checkpoints are not publicly released pending peer-reviewed acceptance, making direct re-use impossible. On a global scale, foundation models like AIFL (Taccari et al., 2026) train multi-million parameter models across 18,588 catchments using the global Caravan dataset (Kratzert et al., 2023). While promising, training and fine-tuning such models requires distributed computing clusters that exceed standard academic workstation budgets.

\section{Flood Early Warning and the National Water Model}

The National Water Model Retrospective version 3.0 provides a 44-year unassimilated physical simulation (February 1979 through January 2023) across CONUS, serving as a standardized physics-based baseline. Investigating model error, Tran et al. (2025) developed Errorcastnet to learn residual corrections for NWM streamflow simulations during extreme events. Their work illustrates how neural networks can operate alongside physical models during flood regimes. In this project, the NWM retrospective simulation serves as an external operational baseline, not an architecture to be corrected or displaced.

\section{Methodological Integrity in Applied Machine Learning}

Experimental validity is a central challenge in data-driven science. Kapoor and Narayanan (2023) surveyed data leakage across 294 studies in 17 scientific fields, finding that temporal lookahead and improper cross-validation frequently produce inflated predictive performance. Their taxonomy directly informs the data partitioning and publication-time filtering implemented here.

To evaluate predictions without unsupported distributional assumptions, I pair each metric with a non-parametric test. For discriminative performance (AUC-ROC and AUC-PR), I use DeLong's test (DeLong et al., 1988) to compare paired model scores on each individual gauge, and Wilcoxon signed-rank tests to compare performance across the panel of test gauges. For warning lead times, I use Kaplan-Meier survival analysis and Cox proportional hazards, treating flood events that fail to trigger an early warning alarm as right-censored at lead time zero. For probabilistic evaluation, I calibrate continuous anomaly scores into event probabilities via logistic regression on validation gauges, and evaluate Brier scores using the Diebold-Mariano test with heteroskedasticity and autocorrelation consistent (HAC) variance estimates (Diebold \& Mariano, 1995).


\chapter{Design and Implementation}

This chapter outlines the system architecture, data pipeline, encoder structure, evaluation framework, and hardware specifications for the proposed investigation. The methodology follows three guiding principles: temporal causality is validated at data ingestion, baseline models receive identical evaluation rigor, and training requirements remain feasible on standard workstation processors.

\section{Proposed System Architecture}

The proposed architecture operates as a sequential pipeline spanning data acquisition, validation, causal feature matrix construction, representation pretraining, anomaly scoring, and statistical evaluation. Figure~\ref{fig:architecture} illustrates the overall information flow through the system.

\begin{figure}[h]
  \centering
  \resizebox{\textwidth}{!}{%
  \begin{tikzpicture}[
    node distance=0.35cm and 0.6cm,
    box/.style={rectangle, draw, rounded corners=2pt, minimum width=1.6cm, minimum height=0.55cm, font=\scriptsize, align=center, thick},
    arrow/.style={-{Stealth[length=2mm]}, semithick}
  ]
    \node[box] (usgs) {USGS\\Streamflow};
    \node[box, below=0.2cm of usgs] (met) {gridMET\\Forcing};
    \node[box, below=0.2cm of met] (snow) {SNODAS\\SWE};
    \node[box, below=0.2cm of snow] (attr) {GAGES-II\\CAMELS};
    \node[box, below=0.2cm of attr] (thresh) {NWPS\\Thresholds};
    \node[box, below=0.2cm of thresh] (events) {NOAA/NSSL\\Events};
    \node[box, below=0.2cm of events] (nwm) {NWM Retro\\v3.0};

    \node[box, right=0.9cm of snow] (valid) {Data\\Validation};
    \node[box, right=0.7cm of valid] (feature) {Causal Feature\\Matrix};
    \node[box, right=0.7cm of feature] (split) {Train/Val/Test\\Split};

    \node[box, above right=0.35cm and 0.7cm of split] (encoder) {SSL\\Encoder};
    \node[box, below=0.2cm of encoder] (scorer) {Anomaly\\Scorer};
    \node[box, below=0.4cm of scorer] (baseline) {4 Baseline\\Categories};

    \node[box, right=0.7cm of scorer] (stats) {Statistical\\Evaluation};

    \foreach \src in {usgs, met, snow, attr, thresh, events, nwm}
      \draw[arrow] (\src) -- (valid);
    \draw[arrow] (valid) -- (feature);
    \draw[arrow] (feature) -- (split);
    \draw[arrow] (split) -- (encoder);
    \draw[arrow] (split) -- (baseline);
    \draw[arrow] (encoder) -- (scorer);
    \draw[arrow] (scorer) -- (stats);
    \draw[arrow] (baseline) -- (stats);
  \end{tikzpicture}%
  }
  \caption{System architecture showing data ingestion from federal repositories, validation checks, causal feature construction, encoder pretraining, anomaly scoring, and multi-baseline statistical evaluation.}
  \label{fig:architecture}
\end{figure}

\subsection{Data Acquisition}

The data ingestion pipeline consists of seven modular collectors, each interacting with an individual federal or scientific repository. Each source has its own collector, so a failure at one provider does not stop the others. All targeted archives represent open public datasets, as summarized in Table~\ref{tab:datasources}.

I select a panel of 54 USGS streamgages distributed across 14 HUC-2 water resource regions in CONUS. The panel is stratified into 27 reference basins (minimal human alteration, identified via GAGES-II) and 27 non-reference basins (regulated flows with upstream dams or diversions), spanning three drainage area tiers: small ($<500\text{ km}^2$), medium ($500\text{--}5000\text{ km}^2$), and large ($>5000\text{ km}^2$). To catch rate limits, connection errors, and schema shifts early, I begin acquisition with an initial pilot of 20 streamgages before expanding to the complete panel.

\begin{table}[h]
  \caption{Data Sources and Variables}
  \label{tab:datasources}
  \centering
  \begin{tabular}{|p{3.2cm}|p{4.5cm}|p{4cm}|}
    \hline
    \textbf{Source} & \textbf{Variables} & \textbf{Access Method} \\
    \hline
    USGS Water Data API & Discharge, gage height & OGC API Features \\
    \hline
    gridMET & Precip., temperature, humidity & THREDDS/OPeNDAP \\
    \hline
    SNODAS (NSIDC) & Snow Water Equivalent & FTP download \\
    \hline
    GAGES-II / CAMELS-US & Static basin attributes & ScienceBase / Zenodo \\
    \hline
    NWPS / CatFIM & Flood-stage thresholds & NWS API \\
    \hline
    NOAA NCEI + NSSL FLASH & Historical flood events & Bulk archive download \\
    \hline
    NWM Retrospective v3.0 & 44-year simulated discharge (1979--2023) & AWS/GCS Zarr \\
    \hline
  \end{tabular}
\end{table}

\subsection{Data Validation and Ground Truth Reconciliation}

Prior to feature extraction, all acquired observations pass through automated quality filters. The validation routines check missingness statistics (discarding gauge records with $>5\%$ missing data across the historical window), verify that observed discharge and stage fall within physically plausible bounds, and confirm temporal continuity.

For evaluation, defining the true occurrence of a flood event requires reconciling multiple sources. I establish an explicit event hierarchy:
\begin{enumerate}
  \item \textbf{Primary Event Truth.} The primary ground truth is defined by observed USGS river stage crossing official NWPS Action stage or higher (Action, Minor, Moderate, or Major flood stage). Action stage marks the threshold where river flows begin approaching bankfull capacity and emergency management surveillance begins.
  \item \textbf{Secondary Corroboration.} Historical records from NOAA NCEI Storm Events and NSSL FLASH serve as secondary corroboration. When an observed stage exceeds Action stage without an NCEI catalog entry, it is retained as an unmodeled local event. If an NCEI storm report lists flooding in the catchment but the gauge stage remains sub-Action, it is flagged as unconfirmed local runoff and excluded from primary lead-time calculations.
\end{enumerate}

\subsection{Causal Feature Matrix and Leakage Prevention}

The feature construction pipeline incorporates three mechanisms designed to prevent temporal leakage and lookahead bias.

\begin{enumerate}
  \item \textbf{Publication-time filter.} A forecast issued at time $t_{\text{issue}}$ may use only values published at or before $t_{\text{issue}}$. For a forecast target $t_{\text{target}}$ evaluated at lead horizon $L$, the issue time is $t_{\text{issue}} = t_{\text{target}} - L$. The loader will store a publication or retrieval timestamp with every observation and drop any observation where $t_{\text{availability}} > t_{\text{issue}}$. For gridMET and SNODAS, where per-value publication timestamps are not available, I use the product's documented near-real-time latency ($\sim$14 hours for gridMET) as an operational proxy and state so where results are reported. The loader logs dropped counts per gauge; an unexpected spike in that count would indicate a timestamp bug.
  \item \textbf{Data vintage tracking.} USGS provisional streamflow observations are subject to retrospective revisions for up to three years before receiving approved status. The pipeline records whether an observation is provisional or approved, tracking the download date and product version for every external input. Because historical revision states cannot be reconstructed for past years in gridMET, documented release latency serves as the operational proxy.
  \item \textbf{Gauge and temporal holdout partitioning.} I partition the 54 streamgages into disjoint subsets: 32 training gauges, 11 validation gauges, and 11 testing gauges, stratified by catchment drainage area and reference status. Within each subset, I enforce a strict chronological cutoff, splitting the multi-decadal record into non-overlapping historical training, validation, and test periods. This two-dimensional split prevents both temporal lookahead and spatial cross-contamination across watersheds (Kapoor \& Narayanan, 2023).
\end{enumerate}

\subsection{Encoder Architecture}

The representation model maps an input tensor $(B, T, C_{\text{in}})$ into a latent sequence $(B, T, D_{\text{latent}})$. I use daily aggregated values rather than the 15-minute instantaneous record. The daily series is longer, quality-controlled, and spans multiple decades; the cost is that rapid flash floods in steep headwater basins, where streams crest within hours, are out of scope.

\begin{itemize}
  \item \textbf{Dilated Causal Convolution Stack.} Convolutions are padded on the left only, so the output at day $t$ depends only on days $\le t$. With kernel size $k=3$ and dilations $d \in \{1, 2, 4, 8\}$, each residual block with two convolutional layers achieves a receptive field of $(k-1) \cdot \sum d = 2 \cdot 15 = 30$ timesteps per block (approximately 31 days). The transformer layers then attend over the entire window, so the model's total historical context is governed by the window length.
  \item \textbf{Causal Transformer Attention.} Pre-LayerNorm transformer blocks apply lower-triangular attention masks so that latent representations at time $t$ attend solely to historical timesteps $\tau \le t$.
  \item \textbf{Causal Masked Reconstruction Objective.} Instead of zeroing values (which is ambiguous because zero discharge is physically valid in ephemeral streams), masked timesteps are replaced by a learned mask token, accompanied by a binary observed/missing indicator channel. Because the convolutional layers use left-sided padding and the transformer layers use causal masks, a masked position at timestep $t$ can attend only to historical observations $\tau \le t$. The pretext task is causal autoregressive span reconstruction—predicting masked future values within the window from preceding context without access to future observations.
\end{itemize}

\begin{table}[h]
  \caption{Encoder Hyperparameters}
  \label{tab:encoder}
  \centering
  \begin{tabular}{|l|l|l|}
    \hline
    \textbf{Parameter} & \textbf{Value} & \textbf{Justification} \\
    \hline
    Latent dimension & 128 & Scaled for multi-channel hydrometry \\
    \hline
    TCN kernel size & 3 & Standard for temporal convolutions \\
    \hline
    Dilation factors & $[1, 2, 4, 8]$ & $\sim$31-day convolutional receptive field \\
    \hline
    Transformer heads / layers & 4 / 2 & Workstation CPU training budget \\
    \hline
    Mask ratio & 15\% & Baseline setting; sensitivity grid: $[10\%, 15\%, 25\%]$ \\
    \hline
    Total parameters & 146,950 & Compact architecture for reproducible CPU training \\
    \hline
  \end{tabular}
\end{table}

\subsection{Anomaly Scoring}

Two scoring formulations are extracted from the pretrained representations to detect emerging hydrological anomalies.

\begin{itemize}
  \item \textbf{Reconstruction Error (Score-A).} During inference, unmasked inputs would yield near-zero reconstruction error for an encoder capable of identity mapping. I therefore specify an inference-time masking protocol: for an evaluation window ending at forecast issue time $t_{\text{issue}}$, the terminal $k=7$ days are replaced with the learned mask token. Score-A is computed as the mean squared error between the model's reconstructed values and the actual observations across these masked days. A high reconstruction error indicates that recent flow dynamics deviate from patterns learned during baseflow regimes.
  \item \textbf{Latent Embedding Distance (Score-B).} The Mahalanobis distance between the terminal latent state and the centroid of baseline training embeddings, capturing distribution shift in representation space.
\end{itemize}

I scale both raw scores per gauge by the median and interquartile range (IQR) computed over an initial reference period, so basins of varying mean discharge can be directly compared.

\subsection{Baseline Methods}

The self-supervised approach is benchmarked against four distinct baseline paradigms.

\begin{enumerate}
  \item \textbf{Physics-Based Simulation.} The National Water Model Retrospective version 3.0 (44-year unassimilated simulation) provides simulated streamflow, evaluated directly against NWPS flood thresholds.
  \item \textbf{Empirical Statistical Methods.} Daily persistence, rolling 30-day day-of-year climatology, and sequential changepoint detectors including cumulative sum (CUSUM) and exponentially weighted moving averages (EWMA).
  \item \textbf{Unsupervised Anomaly Detectors and Ablations.} Isolation Forests and standard LSTM autoencoders trained on the same feature matrices. The LSTM autoencoder isolates the effect of the encoder architecture (TCN-Transformer versus recurrent) while holding the self-supervised reconstruction objective fixed. To further isolate the contribution of self-attention, I also train a TCN-only ablation without transformer layers.
  \item \textbf{Supervised Reference Model.} An Entity-Aware LSTM (Kratzert et al., 2019) trained on flood stage exceedance classifications across six label budgets (1\%, 5\%, 10\%, 25\%, 50\%, 100\% of labeled training windows), evaluated across five fixed random seeds $\{13, 47, 101, 2024, 8891\}$ per budget to establish sample-efficiency curves.
\end{enumerate}

\subsection{Statistical Evaluation}

Each evaluation metric is evaluated using hypothesis testing procedures matched to its distributional properties, as summarized in Table~\ref{tab:stats}.

\begin{table}[h]
  \caption{Statistical Test Assignment}
  \label{tab:stats}
  \centering
  \begin{tabular}{|p{3.0cm}|p{4.2cm}|p{4.8cm}|}
    \hline
    \textbf{Metric Family} & \textbf{Statistical Test} & \textbf{Methodological Rationale} \\
    \hline
    AUC-ROC, AUC-PR & DeLong (within gauge) and Wilcoxon (cross-gauge) & DeLong compares paired models on single gauges; Wilcoxon signed-rank tests across the 11 test gauges; stratified by HUC-2 region \\
    \hline
    Warning Lead Time & Kaplan-Meier and Cox proportional hazards & Evaluates time from first alarm to flood crossing; events with no prior alarm are right-censored at lead 0 \\
    \hline
    Brier Score & Diebold-Mariano with HAC adjustment & Continuous scores calibrated to event probabilities via logistic regression; tests probabilistic sharpness against climatology \\
    \hline
  \end{tabular}
\end{table}

To guard against inflated false positive rates from multiple hypothesis tests, p-values are adjusted using the Benjamini-Hochberg procedure at an overall significance level of $\alpha = 0.05$.

\section{System Requirement Specifications}

\subsection{Software Specifications}

The project dependencies are pinned to ensure reproducible execution across workstations, as listed in Table~\ref{tab:software}.

\begin{table}[h]
  \caption{Key Software Dependencies}
  \label{tab:software}
  \centering
  \begin{tabular}{|l|l|l|}
    \hline
    \textbf{Package} & \textbf{Version} & \textbf{Purpose} \\
    \hline
    Python & 3.12+ & Core runtime \\
    \hline
    PyTorch (CPU) & 2.12.1 & SSL encoder and EA-LSTM baseline \\
    \hline
    scikit-learn & 1.9.0 & Isolation Forest, preprocessing, ROC metrics \\
    \hline
    dataretrieval & 1.3.0 & USGS Water Data API interface \\
    \hline
    lifelines & 0.30.3 & Kaplan-Meier and Cox survival modeling \\
    \hline
    geopandas & 1.1.4 & Catchment geospatial attribute processing \\
    \hline
    xarray & 2026.7.0 & NetCDF and Zarr meteorological arrays \\
    \hline
  \end{tabular}
\end{table}

\subsection{Hardware Specifications}

Table~\ref{tab:hardware} details the hardware requirements for training and evaluating Flood Sentinel.

\begin{table}[h]
  \caption{Hardware Requirements}
  \label{tab:hardware}
  \centering
  \begin{tabular}{|l|l|l|}
    \hline
    \textbf{Component} & \textbf{Minimum} & \textbf{Recommended} \\
    \hline
    CPU & 4 cores & 8+ cores (Apple Silicon or Intel/AMD x86-64) \\
    \hline
    RAM & 8 GB & 16 GB \\
    \hline
    Storage & 20 GB & 50 GB SSD \\
    \hline
    GPU & Not required & Optional for accelerated baseline sweeps \\
    \hline
  \end{tabular}
\end{table}

To confirm computational feasibility on standard workstations, I will conduct an empirical timing pilot on the initial 20-gauge panel using an 8-core workstation CPU. Pretraining for 100 epochs on 20 gauges will measure wall-clock time per epoch and memory footprint, verifying that full-panel training across 54 gauges fits within standard workstation limits.


\chapter{Expected Outcomes and Discussion}

\section{Expected Outcomes}

The project is structured to yield one of three pre-specified findings, with each outcome offering an informative scientific contribution.

\begin{enumerate}
  \item \textbf{Demonstrated Warning Advantage.} Under this positive outcome, the self-supervised anomaly score achieves an AUC-ROC $\ge 0.65$ across the 11 holdout test gauges, with Wilcoxon signed-rank tests confirming statistically significant advantages ($p < 0.05$, Benjamini-Hochberg adjusted) over empirical and unsupervised baselines. Survival analysis reveals positive warning lead times relative to official NWPS flood stages in eligible basins.
  \item \textbf{Inconclusive or Null Discrimination.} If test-panel AUC-ROC falls below 0.55 or the anomaly score provides no advance detection beyond rolling climatology, I will report that as a null result. If AUC-ROC falls between 0.55 and 0.65, the outcome is classified as inconclusive. Across the 11 test gauges, the minimum detectable AUC difference is approximately 0.08 (estimated from pilot score variance at 80\% power), meaning a null result establishes that any effect is smaller than this bound, not that an effect is zero.
  \item \textbf{Sample-Efficiency Crossover Regime.} A third anticipated outcome identifies a sample-efficiency crossover point $k^*$ where the supervised EA-LSTM outperforms the self-supervised anomaly score at high label volumes (e.g., 50\% to 100\% of event labels), while the zero-label self-supervised score matches or exceeds supervised models in label-sparse regimes (e.g., 1\% to 10\% labels). Finding $k^*$ provides empirical guidance on when unsupervised representation learning is preferable to collecting flood labels.
\end{enumerate}

\section{Methodological Highlights}

The methodology incorporates several engineering safeguards designed to ensure scientific transparency and reproducibility:

\begin{itemize}
  \item Five pre-specified research hypotheses with explicit quantitative confirmation, stop, and inconclusive criteria (Appendix~A).
  \item Three-part leakage prevention: publication-time filtering, vintage tracking, and a two-dimensional gauge-and-time split.
  \item Benchmarking against four baseline families: physical simulation (NWM v3.0), empirical filters, unsupervised anomaly detectors, and a supervised EA-LSTM.
  \item Label-budget sweeps across six training-window proportions (1\% to 100\%) evaluated across five fixed random seeds.
  \item Hierarchical ground truth definition prioritizing USGS gage stage crossing NWPS Action stage, with NOAA NCEI and NSSL FLASH catalogs used for secondary verification.
  \item Computational profiling on a 20-gage timing pilot to verify that pretraining converges within standard multi-core workstation budgets.
\end{itemize}

\section{Anticipated Challenges}

Executing the proposed research entails addressing three practical difficulties.

First, the USGS Water Data API undergoes periodic schema and endpoint updates. To catch connection limits, schema shifts, and missing values early, I begin ingestion with an initial pilot of 20 streamgages before pulling the full 54-gauge panel.

Second, severe flood events are sparse in arid and semi-arid basins. With fewer than 5 events in some arid catchments, a Wilcoxon test across those basins alone has low statistical power to detect an AUC difference of 0.05. I therefore stratify evaluation by HUC-2 hydrologic region and treat results in low-prevalence basins as exploratory.

Third, tracking the exact historical release version of gridded meteorological data is technically difficult because archived services like gridMET do not maintain a queryable version history for past decades. I handle this by applying gridMET's documented 14-hour near-real-time publication latency as an operational proxy, filtering observations that postdate forecast issue times.


\chapter{Conclusion and Recommendation}

This proposal fixes the data sources, panel selection, encoder architecture, baselines, and hypotheses before computing test-gauge results. It investigates whether self-supervised pretraining on continuous hydrometric time series produces anomaly scores that give advance warning before streamflow crosses official flood thresholds. Three decisions remain open for resolution during the pilot phase: the exact terminal masking window length $k$ for Score-A, the response-time threshold for lead-time eligibility, and whether revision timestamps can be extracted for a subset of gridMET grids.

\section{Limitations}

Several constraints frame the scope of this planned work:

First, pretraining requires gauges with long, continuous records (at least 20 years of daily data). Consequently, the results will not indicate how well the encoder generalizes to young, sparsely monitored, or ungauged basins, which is where the need for label-free methods is greatest.

Second, I use daily-aggregated discharge and stage records instead of the 15-minute instantaneous series. The daily record is quality-controlled and spans multiple decades; the cost is that rapid flash floods in steep headwater basins, where streams crest within hours, are out of scope. Warning lead times can only be assessed on catchments with response times of 24 hours or longer.

Third, this study is an offline retrospective investigation, not an operational alerting service. It does not account for real-time telemetry delays or sensor dropouts during extreme storms.

\section{Future Enhancement}

Several natural extensions follow the completion of the proposed research:

\begin{itemize}
  \item Evaluating representation transfer to international river basins in the Caravan dataset.
  \item Combining self-supervised latent features with numerical hydrodynamic models to assist physical streamflow simulation.
  \item Incorporating sub-daily radar precipitation feeds to extend precursor detection to smaller headwater catchments.
  \item Incorporating upstream-downstream river network topology using graph neural networks to model hydrologic routing.
\end{itemize}


%% ─── REFERENCES ─────────────────────────────────────────────────────────────
\chapter*{References}
\addcontentsline{toc}{chapter}{References}

\begin{enumerate}[label={[\arabic*]}]
  \item Abatzoglou, J. T. (2013). Development of gridded surface meteorological data for ecological applications. \textit{International Journal of Climatology}, 33(1), 121--131.

  \item Addor, N., Newman, A. J., Mizukami, N., \& Clark, M. P. (2017). The CAMELS data set: Catchment attributes and meteorology for large-sample studies. \textit{Hydrology and Earth System Sciences}, 21(10), 5293--5313.

  \item DeLong, E. R., DeLong, D. M., \& Clarke-Pearson, D. L. (1988). Comparing the areas under two or more correlated receiver operating characteristic curves: A nonparametric approach. \textit{Biometrics}, 44(3), 837--845.

  \item Diebold, F. X., \& Mariano, R. S. (1995). Comparing predictive accuracy. \textit{Journal of Business \& Economic Statistics}, 13(3), 253--263.

  \item Falcone, J. A. (2011). GAGES-II: Geospatial attributes of gages for evaluating streamflow. \textit{U.S. Geological Survey Digital Data Series}.

  \item Haq, S., et al. (2025). HydroGEM: Self-supervised pre-training for streamflow quality control. \textit{arXiv preprint arXiv:2512.14106}.

  \item Kapoor, S., \& Narayanan, A. (2023). Leakage and the reproducibility crisis in machine-learning-based science. \textit{Patterns}, 4(9), 100804.

  \item Kratzert, F., Klotz, D., Shalev, G., Klambauer, G., Hochreiter, S., \& Nearing, G. (2019). Towards learning universal, regional, and local hydrological behaviors via machine learning applied to large-sample datasets. \textit{Hydrology and Earth System Sciences}, 23(12), 5089--5110.

  \item Kratzert, F., et al. (2023). Caravan - A global community dataset for large-sample hydrology. \textit{Scientific Data}, 10, 61.

  \item Liu, C., Liu, D., \& Mu, L. (2022). Improved transformer model for enhanced monthly streamflow predictions of the Yangtze River. \textit{IEEE Access}, 10, 58240--58253.

  \item National Weather Service. (2024). Summary of Natural Hazard Statistics for 2023 in the United States. \textit{NOAA National Weather Service Office of Climate, Water, and Weather Services}.

  \item Nguyen, T.-T.-H., et al. (2023). Streamflow prediction in the Mekong River Basin using deep neural networks. \textit{IEEE Access}, 11, 97930--97943.

  \item Nie, Y., et al. (2023). A time series is worth 64 words: Long-term forecasting with transformers. \textit{ICLR 2023}.

  \item Taccari, M. L., et al. (2026). AIFL: AI foundation model for large-scale streamflow forecasting. \textit{arXiv preprint arXiv:2602.16579}.

  \item Tonekaboni, S., Eytan, D., \& Goldenberg, A. (2021). Unsupervised representation learning for time series with temporal neighborhood coding. \textit{ICLR 2021}.

  \item Tran, Q. Q., et al. (2025). Errorcastnet: Learning error correction for the National Water Model. \textit{AGU Advances}, 6(1), e2024AV001234.

  \item Yue, Z., et al. (2022). TS2Vec: Towards universal representation of time series. \textit{AAAI 2022}, 36(8), 8980--8987.
\end{enumerate}


%% ─── APPENDIX ───────────────────────────────────────────────────────────────
\chapter*{APPENDIX}
\addcontentsline{toc}{chapter}{APPENDIX}
\setcounter{table}{0}
\renewcommand{\thetable}{A.\arabic{table}}

\section*{Appendix A: Pre-Specified Research Hypotheses}

A streamgage catchment is classified as eligible for primary lead-time analysis (H2) if its estimated hydrological response time proxy satisfies $T_{\text{response\_proxy}} \ge 24\text{ h}$ and the gauge has at least one recorded flood stage exceedance during the evaluation period. For headwater catchments ($A \le 500\text{ km}^2$), $T_{\text{response\_proxy}}$ is estimated using the NRCS watershed lag equation; for larger basins ($A > 500\text{ km}^2$), it is estimated from mainstem hydrodynamic reach routing. Basins with $T_{\text{response\_proxy}} < 24\text{ h}$ respond too quickly to be resolved by daily-aggregated records and are excluded from H2.

\begin{table}[h]
  \caption{Research Hypotheses and Falsification Criteria}
  \label{tab:hypotheses}
  \centering
  \begin{tabular}{|p{0.8cm}|p{4.6cm}|p{3.2cm}|p{3.2cm}|}
    \hline
    \textbf{ID} & \textbf{Hypothesis} & \textbf{Evaluation Method} & \textbf{Decision Criteria} \\
    \hline
    H1 & SSL anomaly score achieves test-panel AUC-ROC $\ge$ 0.65 & Wilcoxon signed-rank across 11 test gauges & Confirm: AUC $\ge 0.65$; Stop: AUC $< 0.55$; else Inconclusive \\
    \hline
    H2 & SSL provides $\ge$ 1 day warning lead time for eligible basins & Kaplan-Meier and Cox proportional hazards & Confirm: Median lead $\ge 1$ day; Stop: Median lead $\le 0$ days \\
    \hline
    H3 & SSL discriminates flood events better than NWM Retrospective v3.0 & Wilcoxon signed-rank on paired AUC differences & Confirm: $p < 0.05$ (BH); Stop: $p \ge 0.20$ (BH); else Inconclusive \\
    \hline
    H4 & SSL zero-label anomaly score matches or outperforms EA-LSTM at 100\% labels & Wilcoxon signed-rank on paired AUC differences & Confirm: $p < 0.05$ (BH); Stop: $p \ge 0.20$ (BH); else Inconclusive \\
    \hline
    H5 & Label-budget crossover exists at threshold $k^*$ where SSL beats EA-LSTM & 6 budgets $\times$ 5 random seeds & Confirm: crossover $k^*$ identified; Stop: EA-LSTM beats SSL at all budgets \\
    \hline
  \end{tabular}
\end{table}

\section*{Appendix B: Proposed Development Timeline}

\begin{table}[h]
  \caption{Development Phases and Milestones}
  \label{tab:timeline}
  \centering
  \begin{tabular}{|p{2.2cm}|p{4.8cm}|p{4.8cm}|}
    \hline
    \textbf{Phase} & \textbf{Objective} & \textbf{Key Deliverables} \\
    \hline
    Phase 1 & Pilot data acquisition and pipeline validation & Ingestion across 20 pilot streamgages and initial quality checks \\
    \hline
    Phase 2 & Full panel selection and data extraction & Complete panel of 54 streamgages across 14 HUC-2 regions with basin eligibility estimation \\
    \hline
    Phase 3 & Causal feature engineering and partitioning & Publication-time filtering and 32/11/11 spatial-temporal holdout partitioning \\
    \hline
    Phase 4 & Representation learning and model pretraining & Empirical timing pilot on 8-core workstation and encoder pretraining \\
    \hline
    Phase 5 & Anomaly scoring and baseline execution & Computation of Score-A and Score-B and execution of four baseline families \\
    \hline
    Phase 6 & Statistical hypothesis testing and reporting & DeLong and Wilcoxon tests, Kaplan-Meier lead-time modeling, and final proposal defense \\
    \hline
  \end{tabular}
\end{table}

\end{document}
